% ===== 原文 pp.68--71 =====
% 接续：p.67 末句「Córdova 与 Jafferis 分三步完成了这一约化：」在本块开头续完。

% [原文 p.68]
$\cX(\fg)$ 在 $S^1$ 上约化给出 5d $\cN=2$ SYM；$\cX(\fg)$ 在 $S^3_b$（或其商）上约化
\footnote{$\cX(\fg)$ 在 $S^1\times S^2$ 上的约化也给出复 Chern--Simons，只是能级不同 [311, 312]。}
给出 3d 复 Chern--Simons 理论；$\cX(\fg)$ 在 $S^4_b$ 上约化给出 2d 复 Toda 共形场论。
他们猜测，Toda 共形场论的这种复化版本与通常的 Toda 共形场论对偶。
这一推导在 [313] 中被推广到包含轨形面算子的情形（另见 [314] 的另一种途径）。

\para{约化到复 Chern--Simons 理论。}
约化到 3d 与我们将在 9.3 小节描述的 AGT 对应的 3d/3d 类比有关。
我们把 6d 理论放在 $S^3_b\times C_3$ 上，
其中压扁球面可以用它到 $\R^4$ 的等距嵌入来描述，例如
$S^3_b=\{b^2(y_1^2+y_2^2)+b^{-2}(y_3^2+y_4^2)=r^2\}\subset\R^4$。
保持超对称需要一个部分拓扑扭曲，这相当于给超引力场取合适的背景值，这些值在 [315] 中确定。
[310] 的做法是改用球面 $S^3$ 的另一种压扁方式，
它保持的是 $\mathrm{U}(1)\times\mathrm{SU}(2)$ 等距群而不是 $\mathrm{U}(1)\times\mathrm{U}(1)$。
我们对此略过不提，因为这些背景只相差一些适当的 $Q$-恰当项，最终不影响配分函数。

$S^3$ 的 Hopf 纤维化，即一个以 $S^2$ 为底、以 $S^1$ 为纤维的纤维化，与压扁是相容的。
于是可以先沿 $S^1$ 纤维约化 $\cX(\fg)$，得到 5d $\cN=2$ SYM 理论。
幸亏非阿贝尔 5d 理论有拉氏量描述，从而可以显式地进一步作维数约化；
相比之下，$\cX(\fg)$ 没有已知的拉氏量描述。

沿 Hopf 纤维化的底 $S^2$ 再作一次约化，得到下列轻场，它们都取值于 $\fg$ 的伴随表示。
\begin{itemize}[leftmargin=1.5em,itemsep=2pt,topsep=2pt]
\item 一个 3d 规范场 $A$，来自 5d 规范场沿 $C_3$ 的分量。
它带一个能级 $k=1$ 的 3d Chern--Simons 项，其来源是 5d 引力光子（graviphoton）
在超引力背景中穿过 $S^2$ 有一个单位的通量。这又可追溯到 Hopf 纤维化；
如果改为在 $S^2\times S^1$ 上约化，则 $k=0$。
\item 5d $\cN=2$ SYM 的五个矢量多重态标量的零模。
由于该扭曲把 R 对称性的一个子群 $\fso(3)\subset\fso(5)$ 与 3d 转动等同起来，
这些零模组合成一个一形式 $X$ 和一对标量 $Y_i$。
\item 四个费米子 $\lambda$，其拉氏量含两个导数，为 $-\lambda(\nabla^A)^2\lambda+[X,\lambda]^2$。
用 $\Delta=(\nabla^A)^2+(\mathrm{ad}_X)^2$ 表示，对 $\lambda$ 作二次路径积分给出因子 $(\det\Delta)^2$，
而这对标量 $Y_i$ 因为拉氏量为 $-Y\Delta Y$ 而给出 $1/\det\Delta$。
\end{itemize}

合起来，$Y$ 与 $\lambda$ 贡献因子 $\det\Delta$，
这正好与规范固定 $(\nabla^A)_\mu X^\mu=0$ 的 Faddeev--Popov 行列式相合，
后者对应的“虚”规范变换为 $(A_\mu,X_\mu)\mapsto(A_\mu-[X_\mu,g],X_\mu+\nabla^A_\mu g)$，
其中 $g$ 是局域规范参数。

最终的 3d 理论有一对一形式，从而有一个复一形式 $\mathcal{A}=A+iX\in\fg_\C$，其作用量为
\begin{equation}
S=\frac{q}{8\pi}\int_{C_3}\Tr\Big(\mathcal{A}\wedge\dd\mathcal{A}+\frac{2}{3}\mathcal{A}\wedge\mathcal{A}\wedge\mathcal{A}\Big)
+\frac{\tilde q}{8\pi}\int_{C_3}\Tr\Big(\bar{\mathcal{A}}\wedge\dd\bar{\mathcal{A}}
+\frac{2}{3}\bar{\mathcal{A}}\wedge\bar{\mathcal{A}}\wedge\bar{\mathcal{A}}\Big)
\tag{6.19}
\end{equation}

% [Tutor Note: p.68 三个子弹点的行列式记账可逐项核对：4 个费米子 $\lambda$ 的二次型给出
%   $(\det\Delta)^{4/2}=(\det\Delta)^2$；一对实标量 $Y_i$ 各给 $(\det\Delta)^{-1/2}$，合起来 $1/\det\Delta$；
%   相乘得 $\det\Delta$。而 $(\nabla^A)_\mu X^\mu=0$ 的 Faddeev--Popov 行列式正是 $\det\Delta$（$\Delta=(\nabla^A)^2+(\mathrm{ad}_X)^2$），
%   故虚规范变换恰被这组场抵消，这是作者所说 "matches" 的确切含义。]

% [原文 p.69]
它在来自通常规范不变性与虚规范不变性的 $G_\C$ 规范不变性下成立。
$X$ 的作用量来自许多为保证超对称所必需的超引力场。
耦合常数 $q=k+is$ 与 $\tilde q=k-is$ 以下述方式编码几何
\begin{equation}
k=1,\qquad s=\frac{1-b^2}{1+b^2}.
\tag{6.20}
\end{equation}
更一般地，该理论对 $k\in\Z$ 与 $s\in\R\cup i\R$ 有意义。
[310] 中使用的保持 $\mathrm{SU}(2)\times\mathrm{U}(1)$ 的 $S^3$ 压扁有一个参数 $\ell\in(0,+\infty)$，
且 $s=\sqrt{1-\ell^2}$ 也可以取虚值。
$k$ 的其它取值来自把 $S^3_b$ 换成 $S^2\times S^1$（得 $k=0$），
或换成（压扁的）透镜空间 $L(k,1)_b=S^3_b/\Z_k$。
3d/3d 对应在 9.3 小节讨论。

% [Tutor Note: p.69 式 (6.20) 的 $s\in\R\cup i\R$ 与 $\ell\in(0,+\infty)$ 要一起看：
%   $s=\sqrt{1-\ell^2}$ 在 $\ell<1$ 时为实数、$\ell>1$ 时为纯虚数；而 (6.20) 给 $s=(1-b^2)/(1+b^2)$，
%   由 $|s|<1$ 反推 $b^2=(1-s)/(1+s)>0$，与 $b>0$ 自洽。核对锚点：$b=1\Rightarrow s=0\Rightarrow\ell=1$（未压扁）。]

% [Tutor Note: p.69 把 $b$ 定义为 $\sqrt{\tilde q/q}$，这与 (6.20) 在字面上并不自洽，属文献间的参数化约定差异，
%   正文照原文未改，读者不必以为算错。两种反解：(i) 由 (6.20) 的 $s=(1-b^2)/(1+b^2)$ 解出 $b^2=(1-s)/(1+s)$；
%   (ii) 由 $q=k+is,\ \tilde q=k-is$ 得 $\tilde q/q=(k-is)/(k+is)$，$k=1$ 时即 $(1-is)/(1+is)$，其模恒为 1。
%   对同一个实 $b>0$（$|s|<1$），(ii) 的模为 1 而 (i) 的 $b^2$ 一般不为 1，故二者只在 $b=1$（$s=0$、$\ell=1$）时一致。
%   差别的来源是 $q,\tilde q$ 的定义：若改用 $q=k+s,\ \tilde q=k-s$（即 (6.20) 的 $s$ 不再乘 $i$），
%   则 $\tilde q/q=(1-s)/(1+s)=b^2$ 恰好成立，与 (i) 相符。核对锚点：$b=1\Leftrightarrow s=0\Leftrightarrow q=\tilde q$。]

\para{约化到复 Toda 理论。}
[18] 的想法是把 $S^4_b$ 看成以区间为底、以压扁三球面 $S^3_b$ 为纤维的纤维化。
按 (5.8) 的记号，该区间用 $y_5\in[-r,r]$ 参数化，$S^3_b$ 的半径平方为 $r^2-y_5^2$，
也就是说它在两端都退化为一点。
$S^4_b\times C$ 上的乘积度规 $g$ 经 Weyl 变换映到 $S^3_b\times C_3$ 上的度规，
其中 $C_3$ 是 $C$ 与一个区间的翘曲乘积：
\begin{equation}
g=\dd y_5^2+(r^2-y_5^2)\,g_{S^3_b}+g_C,\qquad
\frac{1}{r^2-y_5^2}\,g=g_{S^3_b}+\frac{\dd y_5^2+g_C}{r^2-y_5^2}.
\tag{6.21}
\end{equation}
所得度规在 $y_5=\pm r$ 处奇异，这给 $C_3=[-r,r]\times_w C$ 上的理论带来边界条件。
来自两端的边缘模于是被理解为由手征复 Toda 理论描述。
把这两个手征理论合起来，就得到 $C$ 上的复 Toda 共形场论。
该理论描述 Cartan 子代数 $\fh\subset\fg$ 的复化中的一个复玻色子 $\Phi$，带指数型势：
\begin{equation}
S_{\fg_\C\ \text{Toda}}
=\frac{q}{8\pi}\int_C\Big(\langle\partial\Phi,\bar\partial\Phi\rangle
+\sum_{j=1}^{r}\ee^{\langle e_j,\Phi\rangle}\Big)\dd^2z
+\frac{\tilde q}{8\pi}\int_C\Big(\langle\partial\bar\Phi,\bar\partial\bar\Phi\rangle
+\sum_{j=1}^{r}\ee^{\langle e_j,\bar\Phi\rangle}\Big)\dd^2z,
\tag{6.22}
\end{equation}
其中 $e_j$ 是 $\fg$ 的单根，$r=\rank\fg$，Killing 型与 $\fh^*$、$\fh$ 之间的配对都记作 $\langle\ ,\ \rangle$。
耦合常数 $q=k+is$、$\tilde q=k-is$ 按 (6.20) 编码几何。

\para{与通常 Toda 理论的关系。}
接下来的猜想是：复 Toda 共形场论与 [22] 早先针对 $S^4_b/\Z_k$ 上的 AGT 对应所提的设想有关
（其中 $k=2$ 的情形基于 [20]）。
对 $\fg=\fsu(N)$，[22] 所提的 2d CFT 由两个退耦的理论组成：
一个 $\widehat{\fsu}(k)_N/\widehat{\fu}(1)^{k-1}$ 陪集，以及参数为 $N$、$k$
与 $b=\sqrt{\tilde q/q}$（与压扁参数吻合）的实 parafermion Toda 共形场论。
后一理论描述 parafermion 与实玻色子 $\psi,\phi\in\fh$，
其中 parafermion 由另一个陪集模型 $\widehat{\fsu}(N)_k/\widehat{\fu}(1)^{N-1}$ 描述，
并通过维数为 $1-1/k$ 的算子 $\psi_j\bar\psi_j$ 与实玻色子耦合，
\begin{equation}
S_{\text{para-Toda}}
=S\Big(\frac{\widehat{\fsu}(N)_k}{\widehat{\fu}(1)^{N-1}}\Big)
+\int_C\Big(\langle\partial\phi,\bar\partial\phi\rangle
+\sum_{j=1}^{N-1}\psi_j\bar\psi_j\,\ee^{(b/\sqrt{k})\langle e_j,\phi\rangle}\Big)\dd^2z.
\tag{6.23}
\end{equation}

% [Tutor Note: p.69 式 (6.22) 的横杠位置必须逐符号核对：第一项是 $\langle\partial\Phi,\bar\partial\Phi\rangle$（只有第二个 $\partial$ 带杠），
%   第二项是 $\langle\partial\bar\Phi,\bar\partial\bar\Phi\rangle$ 与 $\ee^{\langle e_j,\bar\Phi\rangle}$（$\Phi$ 全带杠）。
%   漏一个杠就会破坏 $q$ 与 $\tilde q$ 两部分的复共轭结构；核对办法：$k=1$ 且 $s=0$（即 (6.20) 给出 $b=1$）时 $q=\tilde q=1$，
%   此时 (6.22) 的两项互为复共轭，相加给出实 Toda 作用量的两倍，这正对应正文所说「$k=1$ 时只剩耦合为 $b$ 的通常 Toda」。]

% [原文 p.70]
当 $k=1$ 时，退耦的陪集与 parafermion 都变得平凡，只剩下耦合为 $b$ 的通常 Toda 共形场论，
这正是标准 AGT 对应所断言的内容。
复 Toda 共形场论与「陪集加 para-Toda 共形场论」之间的猜想对偶，
只在最简单的情形下被显式检验过：$\fg=\fsu(2)$ 且 $k=1$ 的情形见 [316]，
$N=k=2$ 的情形见 [29, 317]，后者本质上归结为玻色化。
关于轨形情形的进一步讨论见第 83 页。

\part{第三篇\quad AGT 的推广}

\section{一般 class S 理论}

本节我们进入非拉氏量理论的领域：虽然所有由 $\cX(\fsu(2))$ 加驯服穿刺得到的 class S 理论
都可以通过耦合矢量多重态与超多重态来实现，我们现在从两个方面推广这一图景。

6d $(2,0)$ 理论 $\cX(\fg)$ 由任意的单链单李代数 $\fg$ 标记，
因此 AGT 对应 [5] 从 $\fsu(2)$ 推广到 $\fsu(N)$ [93] 以及一般的规范代数 [96, 98, 99, 318]、
Liouville 共形场论推广为 Toda 共形场论，也就不足为奇了。
在 7.1 小节中带着「与规范理论的联系」这一着眼点回顾该共形场论之后，
我们在 7.2 小节描述 AGT 对应的一个例子，以及关于穿刺点的一些重要考量。

% 原文 p.70 这个要点框并无标题条（与 §1--§6 各节的框不同），故标题参数留空。
\begin{keyboxplain}
\begin{itemize}[leftmargin=1.4em,itemsep=1pt,topsep=1pt]
\item 大多数高阶 class S 理论是非拉氏量的。
\item 部分 Higgs 化给出驯服穿刺点的层级结构。其中一些由 SU 群的箭图尾巴描述。
\item 驯服穿刺点相撞给出野生穿刺点。这往往导致 AD 理论。
\item 驯服穿刺点映到 Toda 半退化主算子。
\item 野生穿刺点映到 Toda 共形场论的非正则算子。
\end{itemize}
\end{keyboxplain}

其次，在 7.3 小节中，我们考虑这样的有趣极限：两个穿刺点相撞，
而描述这些缺陷的参数作适当的标度。
由此得到的野生穿刺点可以给出渐近自由的规范理论（例如 $N_f<4$ 的 $\mathrm{SU}(2)$ SQCD），
以及作为 class S 一部分的 AD 理论 [2]。

% [Tutor Note: p.70 要点框在原排版中是右缘绕排的浮动框，正文被压成一半宽度（这是原图的版式，
%   不是「本节重点」之类的编者加注）。这里按常规单栏排版，框放在第二段之后，
%   因为原图中「Second, in subsection 7.3…」一段正是绕框而行、接在「…about punctures in subsection 7.2.」之后的。
%   核对方法：p.70.png 左侧前五行行宽明显短于正文栏宽，且右边界恰好贴着灰框左边。]

\subsection{Toda 共形场论}

\para{拉氏量与对称性。}
我们无法公正对待 Toda 理论五十年的历史，从 1967 年的 Toda 晶格 [319, 320] 讲起------
那是由带最近邻指数相互作用的质点组成的系统（早期综述见 [321]）。
它的量子场论版本于 1982 年在 [322] 中定义，并被发现是共形的。
给定单链
\footnote{Toda 共形场论对任意单李代数都有定义，但我们只给出单链情形，
因为只有这一情形与 AGT 对应相关。
非单链 Toda 共形场论中的关联函数，想必对应于 $S^4_b$ 上规范理论中适当的外自同构扭曲算子的关联函数。}
李代数 $\fg$ 及其 Cartan 李代数 $\fh$，Toda 共形场论描述一个实标量场 $\phi\in\fh$，

% [原文 p.71]
其拉氏量密度为
\begin{equation}
\mathcal{L}=\frac{1}{8\pi}\Big(\hat g^{ab}\langle\partial_a\phi,\partial_b\phi\rangle
+2\langle Q,\phi\rangle\hat R\Big)
+\mu\sum_{i=1}^{\rank\fg}\ee^{b\langle e_i,\phi\rangle}.
\tag{7.1}
\end{equation}
这里 $e_i$ 是 $\fg$ 的单根，$\langle\ ,\ \rangle$ 既表示 $\fh^*$ 与 $\fh$ 的配对，也表示 Killing 型。
乘在背景度规 $\hat g$ 的标量曲率 $\hat R$ 上的背景荷矢量 $Q$ 取为 $Q=(b+\frac{1}{b})\rho$，
其中 $\rho$ 是 Weyl 矢量，即正根之和的一半，等价地，是全部基本权 $\varpi_j$ 之和。
顶点算子 $V_\alpha=\ee^{\langle\alpha,\phi\rangle}$（在这一记号中我们略去正规序）
由 $\alpha\in\fh^*_\C$ 标记，其全纯共形维数为
\footnote{Liouville 与 Toda 共形场论中 $Q$ 的标准约定相差一个因子 2，
因此典型的 Liouville 动量形如 $\alpha=Q/2+iP$，而典型的 Toda 动量是 $\alpha=Q+iP$。}
\begin{equation}
h(\alpha)=\frac{1}{2}\langle\alpha,2Q-\alpha\rangle.
\tag{7.2}
\end{equation}
特别地 $h(be_i)=1$，这保证了指数型势项恰好是边缘的。
它们的耦合 $\mu$ 是冗余的，其效果只是把 $\phi$ 平移 $\rho$ 的某个倍数。
但耦合 $b>0$ 起本质作用：例如中心荷 $c=\rank\fg+12\langle Q,Q\rangle$ 就依赖于它。
对单链 $\fg$，该理论（被认为）在 $b\mapsto 1/b$ 下对偶，同时保持
$\lambda=[\pi\Gamma(b^2)\mu/\Gamma(1-b^2)]^{1/b}$ 不变。
这对 AGT 对应来说是相当令人满意的，因为 $S^4_b$ 与 $S^4_{1/b}$ 是等距的。

除 2d CFT 的无穷维 Virasoro 对称性之外，Toda 共形场论还有（反）全纯的 $W_\fg$ 对称性。
这一手征代数在 [323] 中对 $\fg=\fsu(3)$ 被发现，
更广泛地，在 [324--326] 中作为极小模型的对称性被发现。早期综述见 [327]。
\footnote{同一批作者的同时期文章 [328] 据说与此相关，但我没能拿到。}
它可以通过仿射李代数的量子 Drinfeld--Sokolov 约化来实现 [131, 132, 329, 330]。
（在 W-弦方面的应用见 [331, 332]。）
在 $\fg=\fsu(N)$ 情形，它可以实现为由无穷多个守恒流生成的更一般手征代数 $W_\infty$ 的截断，
综述见 [134]。作为手征代数，$W_\fg$ 由 $\rank\fg$ 个守恒流 $W^{(k)}(z)$ 生成，
这些流的自旋 $k$ 是 $\fg$ 的 Casimir 不变量的次数。
二次 Casimir 不变量给出全纯能动张量 $T(z)=W^{(2)}(z)$，它生成 $W_\fg$ 的 Virasoro 子代数。

\para{主算子与归一化。}
顶点算子 $V_\alpha=\ee^{\langle\alpha,\phi\rangle}$ 在所有 $W^{(k)}$ 的零模下有确定的量子数 $w^{(k)}(\alpha)$，
也就是说 OPE 以如下方式开始
\begin{equation}
W^{(k)}(z)V_\alpha(0)=\frac{w^{(k)}(\alpha)}{z^k}V_\alpha(0)+\cdots
\tag{7.3}
\end{equation}
例如 $w^{(2)}(\alpha)=h(\alpha)$ 由 (7.2) 给出。
守恒流 $W^{(k)}(z)$ 在规范理论一侧对应于 $k$ 次微分 $\phi_k$，
后者出现在 Seiberg--Witten 曲线的构造 (1.4) 中。
动量 $\Im\alpha$ 是某个给定驯服穿刺点处（对角化的）质量参数 $m\in\fg$ 的 $r$ 倍。
在经典极限 $r\to\infty$ 下，量子的

% [Tutor Note: p.71 式 (7.2) 有一个可逐行验算的锚点，即正文紧接的 $h(be_i)=1$：
%   取单链标准归一化 $\langle e_i,e_i\rangle=2$，则 $\langle e_i,\varpi_j\rangle=\delta_{ij}$，
%   而 $\rho=\sum_j\varpi_j$ 给出 $\langle e_i,\rho\rangle=1$。于是由 $Q=(b+\frac1b)\rho$ 得 $\langle e_i,Q\rangle=b+\frac1b$，
%   代入 $h(be_i)=\frac12\langle be_i,2Q-be_i\rangle=b\langle e_i,Q\rangle-\frac{b^2}{2}\langle e_i,e_i\rangle
%   =b(b+\frac1b)-b^2=1$。这正说明 (7.2) 的 $\frac12$ 因子与 $\langle e_i,e_i\rangle=2$ 的归一化是配套的，
%   也解释了为何指数型势项 $\ee^{b\langle e_i,\phi\rangle}$ 在 $h=1$ 时恰好边缘。]
